% Options for packages loaded elsewhere
% Options for packages loaded elsewhere
\PassOptionsToPackage{unicode}{hyperref}
\PassOptionsToPackage{hyphens}{url}
\PassOptionsToPackage{dvipsnames,svgnames,x11names}{xcolor}
%
\documentclass[
  letterpaper,
  DIV=11,
  numbers=noendperiod]{scrartcl}
\usepackage{xcolor}
\usepackage{amsmath,amssymb}
\setcounter{secnumdepth}{-\maxdimen} % remove section numbering
\usepackage{iftex}
\ifPDFTeX
  \usepackage[T1]{fontenc}
  \usepackage[utf8]{inputenc}
  \usepackage{textcomp} % provide euro and other symbols
\else % if luatex or xetex
  \usepackage{unicode-math} % this also loads fontspec
  \defaultfontfeatures{Scale=MatchLowercase}
  \defaultfontfeatures[\rmfamily]{Ligatures=TeX,Scale=1}
\fi
\usepackage{lmodern}
\ifPDFTeX\else
  % xetex/luatex font selection
\fi
% Use upquote if available, for straight quotes in verbatim environments
\IfFileExists{upquote.sty}{\usepackage{upquote}}{}
\IfFileExists{microtype.sty}{% use microtype if available
  \usepackage[]{microtype}
  \UseMicrotypeSet[protrusion]{basicmath} % disable protrusion for tt fonts
}{}
\makeatletter
\@ifundefined{KOMAClassName}{% if non-KOMA class
  \IfFileExists{parskip.sty}{%
    \usepackage{parskip}
  }{% else
    \setlength{\parindent}{0pt}
    \setlength{\parskip}{6pt plus 2pt minus 1pt}}
}{% if KOMA class
  \KOMAoptions{parskip=half}}
\makeatother
% Make \paragraph and \subparagraph free-standing
\makeatletter
\ifx\paragraph\undefined\else
  \let\oldparagraph\paragraph
  \renewcommand{\paragraph}{
    \@ifstar
      \xxxParagraphStar
      \xxxParagraphNoStar
  }
  \newcommand{\xxxParagraphStar}[1]{\oldparagraph*{#1}\mbox{}}
  \newcommand{\xxxParagraphNoStar}[1]{\oldparagraph{#1}\mbox{}}
\fi
\ifx\subparagraph\undefined\else
  \let\oldsubparagraph\subparagraph
  \renewcommand{\subparagraph}{
    \@ifstar
      \xxxSubParagraphStar
      \xxxSubParagraphNoStar
  }
  \newcommand{\xxxSubParagraphStar}[1]{\oldsubparagraph*{#1}\mbox{}}
  \newcommand{\xxxSubParagraphNoStar}[1]{\oldsubparagraph{#1}\mbox{}}
\fi
\makeatother

\usepackage{color}
\usepackage{fancyvrb}
\newcommand{\VerbBar}{|}
\newcommand{\VERB}{\Verb[commandchars=\\\{\}]}
\DefineVerbatimEnvironment{Highlighting}{Verbatim}{commandchars=\\\{\}}
% Add ',fontsize=\small' for more characters per line
\usepackage{framed}
\definecolor{shadecolor}{RGB}{241,243,245}
\newenvironment{Shaded}{\begin{snugshade}}{\end{snugshade}}
\newcommand{\AlertTok}[1]{\textcolor[rgb]{0.68,0.00,0.00}{#1}}
\newcommand{\AnnotationTok}[1]{\textcolor[rgb]{0.37,0.37,0.37}{#1}}
\newcommand{\AttributeTok}[1]{\textcolor[rgb]{0.40,0.45,0.13}{#1}}
\newcommand{\BaseNTok}[1]{\textcolor[rgb]{0.68,0.00,0.00}{#1}}
\newcommand{\BuiltInTok}[1]{\textcolor[rgb]{0.00,0.23,0.31}{#1}}
\newcommand{\CharTok}[1]{\textcolor[rgb]{0.13,0.47,0.30}{#1}}
\newcommand{\CommentTok}[1]{\textcolor[rgb]{0.37,0.37,0.37}{#1}}
\newcommand{\CommentVarTok}[1]{\textcolor[rgb]{0.37,0.37,0.37}{\textit{#1}}}
\newcommand{\ConstantTok}[1]{\textcolor[rgb]{0.56,0.35,0.01}{#1}}
\newcommand{\ControlFlowTok}[1]{\textcolor[rgb]{0.00,0.23,0.31}{\textbf{#1}}}
\newcommand{\DataTypeTok}[1]{\textcolor[rgb]{0.68,0.00,0.00}{#1}}
\newcommand{\DecValTok}[1]{\textcolor[rgb]{0.68,0.00,0.00}{#1}}
\newcommand{\DocumentationTok}[1]{\textcolor[rgb]{0.37,0.37,0.37}{\textit{#1}}}
\newcommand{\ErrorTok}[1]{\textcolor[rgb]{0.68,0.00,0.00}{#1}}
\newcommand{\ExtensionTok}[1]{\textcolor[rgb]{0.00,0.23,0.31}{#1}}
\newcommand{\FloatTok}[1]{\textcolor[rgb]{0.68,0.00,0.00}{#1}}
\newcommand{\FunctionTok}[1]{\textcolor[rgb]{0.28,0.35,0.67}{#1}}
\newcommand{\ImportTok}[1]{\textcolor[rgb]{0.00,0.46,0.62}{#1}}
\newcommand{\InformationTok}[1]{\textcolor[rgb]{0.37,0.37,0.37}{#1}}
\newcommand{\KeywordTok}[1]{\textcolor[rgb]{0.00,0.23,0.31}{\textbf{#1}}}
\newcommand{\NormalTok}[1]{\textcolor[rgb]{0.00,0.23,0.31}{#1}}
\newcommand{\OperatorTok}[1]{\textcolor[rgb]{0.37,0.37,0.37}{#1}}
\newcommand{\OtherTok}[1]{\textcolor[rgb]{0.00,0.23,0.31}{#1}}
\newcommand{\PreprocessorTok}[1]{\textcolor[rgb]{0.68,0.00,0.00}{#1}}
\newcommand{\RegionMarkerTok}[1]{\textcolor[rgb]{0.00,0.23,0.31}{#1}}
\newcommand{\SpecialCharTok}[1]{\textcolor[rgb]{0.37,0.37,0.37}{#1}}
\newcommand{\SpecialStringTok}[1]{\textcolor[rgb]{0.13,0.47,0.30}{#1}}
\newcommand{\StringTok}[1]{\textcolor[rgb]{0.13,0.47,0.30}{#1}}
\newcommand{\VariableTok}[1]{\textcolor[rgb]{0.07,0.07,0.07}{#1}}
\newcommand{\VerbatimStringTok}[1]{\textcolor[rgb]{0.13,0.47,0.30}{#1}}
\newcommand{\WarningTok}[1]{\textcolor[rgb]{0.37,0.37,0.37}{\textit{#1}}}

\usepackage{longtable,booktabs,array}
\usepackage{calc} % for calculating minipage widths
% Correct order of tables after \paragraph or \subparagraph
\usepackage{etoolbox}
\makeatletter
\patchcmd\longtable{\par}{\if@noskipsec\mbox{}\fi\par}{}{}
\makeatother
% Allow footnotes in longtable head/foot
\IfFileExists{footnotehyper.sty}{\usepackage{footnotehyper}}{\usepackage{footnote}}
\makesavenoteenv{longtable}
\usepackage{graphicx}
\makeatletter
\newsavebox\pandoc@box
\newcommand*\pandocbounded[1]{% scales image to fit in text height/width
  \sbox\pandoc@box{#1}%
  \Gscale@div\@tempa{\textheight}{\dimexpr\ht\pandoc@box+\dp\pandoc@box\relax}%
  \Gscale@div\@tempb{\linewidth}{\wd\pandoc@box}%
  \ifdim\@tempb\p@<\@tempa\p@\let\@tempa\@tempb\fi% select the smaller of both
  \ifdim\@tempa\p@<\p@\scalebox{\@tempa}{\usebox\pandoc@box}%
  \else\usebox{\pandoc@box}%
  \fi%
}
% Set default figure placement to htbp
\def\fps@figure{htbp}
\makeatother





\setlength{\emergencystretch}{3em} % prevent overfull lines

\providecommand{\tightlist}{%
  \setlength{\itemsep}{0pt}\setlength{\parskip}{0pt}}



 


\KOMAoption{captions}{tableheading}
\makeatletter
\@ifpackageloaded{caption}{}{\usepackage{caption}}
\AtBeginDocument{%
\ifdefined\contentsname
  \renewcommand*\contentsname{Table of contents}
\else
  \newcommand\contentsname{Table of contents}
\fi
\ifdefined\listfigurename
  \renewcommand*\listfigurename{List of Figures}
\else
  \newcommand\listfigurename{List of Figures}
\fi
\ifdefined\listtablename
  \renewcommand*\listtablename{List of Tables}
\else
  \newcommand\listtablename{List of Tables}
\fi
\ifdefined\figurename
  \renewcommand*\figurename{Figure}
\else
  \newcommand\figurename{Figure}
\fi
\ifdefined\tablename
  \renewcommand*\tablename{Table}
\else
  \newcommand\tablename{Table}
\fi
}
\@ifpackageloaded{float}{}{\usepackage{float}}
\floatstyle{ruled}
\@ifundefined{c@chapter}{\newfloat{codelisting}{h}{lop}}{\newfloat{codelisting}{h}{lop}[chapter]}
\floatname{codelisting}{Listing}
\newcommand*\listoflistings{\listof{codelisting}{List of Listings}}
\makeatother
\makeatletter
\makeatother
\makeatletter
\@ifpackageloaded{caption}{}{\usepackage{caption}}
\@ifpackageloaded{subcaption}{}{\usepackage{subcaption}}
\makeatother
\usepackage{bookmark}
\IfFileExists{xurl.sty}{\usepackage{xurl}}{} % add URL line breaks if available
\urlstyle{same}
\hypersetup{
  pdftitle={One Hundredth of a Second},
  pdfauthor={Kristian Vepsäläinen},
  colorlinks=true,
  linkcolor={blue},
  filecolor={Maroon},
  citecolor={Blue},
  urlcolor={Blue},
  pdfcreator={LaTeX via pandoc}}


\title{One Hundredth of a Second}
\usepackage{etoolbox}
\makeatletter
\providecommand{\subtitle}[1]{% add subtitle to \maketitle
  \apptocmd{\@title}{\par {\large #1 \par}}{}{}
}
\makeatother
\subtitle{A Bayesian forward simulation of aerodynamic drag at the 1980
Olympics}
\author{Kristian Vepsäläinen}
\date{2026-07-01}
\begin{document}
\maketitle


\subsection{}\label{section}

One hundredthof a second.

Lake Placid, 17 February 1980

Kristian Vepsäläinen ~·~ useR! 2026

\begin{center}\rule{0.5\linewidth}{0.5pt}\end{center}

\subsection{The race}\label{the-race}

\textbf{Men's 15 km cross-country skiing}

\begin{itemize}
\tightlist
\item
  Gold: \textbf{Thomas Wassberg} ~ 41:57.63
\item
  Silver: \textbf{Juha Mieto} ~~~~~~~ 41:57.64
\end{itemize}

Δ = 0.01 s

After 42 minutes of skiing, one hundredth of a second. The IOC changed
the rules. Times are now rounded to 0.1 s.

\textbf{The folk theory}

Juha Mieto had a beard. A large, full beard.

For decades Finnish skiing folklore has asked:

\emph{Did the beard cost him the gold?}

\begin{center}\rule{0.5\linewidth}{0.5pt}\end{center}

\subsection{This is not a history
talk}\label{this-is-not-a-history-talk}

We \textbf{cannot} know what happened on 17 February 1980.

We \textbf{can} ask a different, sharper question:

\begin{quote}
Under physically plausible assumptions, is \textbf{0.01 s} a
\emph{typical} outcome of additional aerodynamic drag, or an
\emph{extreme} one?
\end{quote}

This is a \textbf{modeling} question --- and it is exactly the kind of
question R is good at.

\begin{center}\rule{0.5\linewidth}{0.5pt}\end{center}

\subsection{Why this is a useR talk}\label{why-this-is-a-user-talk}

This is a story about three things every applied modeller meets:

\begin{enumerate}
\def\labelenumi{\arabic{enumi}.}
\tightlist
\item
  \textbf{Sparse data.} We have three split times. That's it.
\item
  \textbf{Unknown nuisance parameters.} Course profile, wind, body
  shape, power output.
\item
  \textbf{A claim about a small margin.} ``Negligible''
  vs.~``plausible'' --- who decides?
\end{enumerate}

The toolkit:

\textbf{Bayesian forward simulation} + \textbf{R} + \textbf{honest
priors} + \textbf{a single plot}.

\begin{center}\rule{0.5\linewidth}{0.5pt}\end{center}

\subsection{The plan}\label{the-plan}

\subsubsection{1. Physics}\label{physics}

A drag-power model of cross-country skiing.

Frontal area · drag coefficient · velocity³ · power.

\subsubsection{2. Priors}\label{priors}

Honest, wide, defensible ranges for every unknown.

Not point estimates.

\subsubsection{3. Propagate}\label{propagate}

100,000 simulated races.

Look at the posterior of the time loss attributable to drag.

No Stan. No JAGS. Just \texttt{runif}, \texttt{rnorm}, and a loop.

\begin{center}\rule{0.5\linewidth}{0.5pt}\end{center}

\subsection{What we know}\label{what-we-know}

\textbf{Three split times for Juha Mieto} (sports-reference.com archive)

\begin{longtable}[]{@{}rrr@{}}
\toprule\noalign{}
Distance & Cumulative time & Split speed \\
\midrule\noalign{}
\endhead
\bottomrule\noalign{}
\endlastfoot
5 km & 13:25.00 & 6.21 m/s \\
10 km & 25:51.97 & 6.70 m/s \\
15 km & 41:57.64 & 5.18 m/s \\
\end{longtable}

\textbf{Venue:} Mt. Van Hoevenberg, Lake Placid, NY. The trails still
exist.

\textbf{Everything else is a prior.}

\begin{center}\rule{0.5\linewidth}{0.5pt}\end{center}

\subsection{The physics}\label{the-physics}

The power needed to overcome aerodynamic drag:

\[
P_\text{drag} \;=\; \tfrac{1}{2}\,\rho\,C_d\,A\,v^{3}
\]

\begin{longtable}[]{@{}
  >{\raggedright\arraybackslash}p{(\linewidth - 4\tabcolsep) * \real{0.2857}}
  >{\raggedright\arraybackslash}p{(\linewidth - 4\tabcolsep) * \real{0.3214}}
  >{\raggedright\arraybackslash}p{(\linewidth - 4\tabcolsep) * \real{0.3929}}@{}}
\toprule\noalign{}
\begin{minipage}[b]{\linewidth}\raggedright
Symbol
\end{minipage} & \begin{minipage}[b]{\linewidth}\raggedright
Meaning
\end{minipage} & \begin{minipage}[b]{\linewidth}\raggedright
Treatment
\end{minipage} \\
\midrule\noalign{}
\endhead
\bottomrule\noalign{}
\endlastfoot
\(\rho\) & Air density & Fixed at 1.2 kg/m³ \\
\(C_d\) & Drag coefficient (body) & Prior: U(0.7, 0.9) \\
\(A\) & Frontal area (body) & Fixed at 0.5 m² \\
\(A_\text{beard}\) & Additional frontal area from beard & Prior:
U(0.005, 0.015) m² \\
\(v\) & Local velocity & Derived from splits + terrain \\
\(P\) & Total power output & Prior: N(400, 30) W \\
\end{longtable}

The beard adds frontal area. The body still has to push through air.
\(\Delta P = \tfrac{1}{2}\rho C_d A_\text{beard} v^3 \cdot \alpha^{3}\)
where \(\alpha < 1\) accounts for slipstream shielding.

\begin{center}\rule{0.5\linewidth}{0.5pt}\end{center}

\subsection{Anchoring v to the split
times}\label{anchoring-v-to-the-split-times}

The splits give us \textbf{average segment speeds}. We don't know how
the speed was distributed within each segment.

So we model it:

\begin{Shaded}
\begin{Highlighting}[]
\NormalTok{profile     }\OtherTok{\textless{}{-}} \FunctionTok{c}\NormalTok{(}\AttributeTok{up =} \FloatTok{0.4}\NormalTok{, }\AttributeTok{flat =} \FloatTok{0.3}\NormalTok{, }\AttributeTok{down =} \FloatTok{0.3}\NormalTok{)   }\CommentTok{\# course mix}
\NormalTok{speed\_mult  }\OtherTok{\textless{}{-}} \FunctionTok{c}\NormalTok{(}\AttributeTok{up =} \FloatTok{0.7}\NormalTok{, }\AttributeTok{flat =} \FloatTok{1.0}\NormalTok{, }\AttributeTok{down =} \FloatTok{1.3}\NormalTok{)   }\CommentTok{\# relative speeds}

\ControlFlowTok{for}\NormalTok{ (k }\ControlFlowTok{in} \DecValTok{1}\SpecialCharTok{:}\DecValTok{3}\NormalTok{) \{                                      }\CommentTok{\# for each split}
\NormalTok{  v\_obs }\OtherTok{\textless{}{-}}\NormalTok{ v\_segments[k]}
\NormalTok{  scale }\OtherTok{\textless{}{-}}\NormalTok{ v\_obs }\SpecialCharTok{/} \FunctionTok{sum}\NormalTok{(profile }\SpecialCharTok{*}\NormalTok{ speed\_mult)          }\CommentTok{\# match the split}
  \ControlFlowTok{for}\NormalTok{ (seg }\ControlFlowTok{in} \FunctionTok{names}\NormalTok{(profile)) \{}
\NormalTok{    v }\OtherTok{\textless{}{-}} \FunctionTok{rnorm}\NormalTok{(N, speed\_mult[seg] }\SpecialCharTok{*}\NormalTok{ scale, }\FloatTok{0.2}\NormalTok{)       }\CommentTok{\# local v}
\NormalTok{    v[v }\SpecialCharTok{\textless{}} \DecValTok{1}\NormalTok{] }\OtherTok{\textless{}{-}} \DecValTok{1}
\NormalTok{    v3\_mean }\OtherTok{\textless{}{-}}\NormalTok{ v3\_mean }\SpecialCharTok{+}\NormalTok{ profile[seg] }\SpecialCharTok{*}\NormalTok{ v}\SpecialCharTok{\^{}}\DecValTok{3}           \CommentTok{\# accumulate E[v³]}
\NormalTok{  \}}
\NormalTok{\}}
\end{Highlighting}
\end{Shaded}

Jensen's inequality matters: drag depends on \(\mathbb{E}[v^3]\), not
\(\mathbb{E}[v]^3\). Modelling within-segment variation is \textbf{not
optional}.

\begin{center}\rule{0.5\linewidth}{0.5pt}\end{center}

\subsection{The priors, visualized}\label{the-priors-visualized}

\pandocbounded{\includegraphics[keepaspectratio]{mieto-wassberg-useR2026_files/figure-pdf/prior-viz-1.png}}

Wide on purpose. The goal is to ask: \emph{does the conclusion survive
honest uncertainty?} --- not to pretend we have a tight estimate.

\begin{center}\rule{0.5\linewidth}{0.5pt}\end{center}

\subsection{The model --- one screen of
R}\label{the-model-one-screen-of-r}

\begin{Shaded}
\begin{Highlighting}[]
\FunctionTok{set.seed}\NormalTok{(}\DecValTok{1980}\NormalTok{)}
\NormalTok{N }\OtherTok{\textless{}{-}} \DecValTok{100000}\NormalTok{;  rho }\OtherTok{\textless{}{-}} \FloatTok{1.2}\NormalTok{;  A\_body }\OtherTok{\textless{}{-}} \FloatTok{0.5}
\NormalTok{T\_total }\OtherTok{\textless{}{-}} \DecValTok{41}\SpecialCharTok{*}\DecValTok{60} \SpecialCharTok{+} \FloatTok{57.64}

\NormalTok{v\_segments }\OtherTok{\textless{}{-}} \FunctionTok{c}\NormalTok{(}\DecValTok{5000} \SpecialCharTok{/}\NormalTok{ (}\DecValTok{13}\SpecialCharTok{*}\DecValTok{60} \SpecialCharTok{+} \FloatTok{25.00}\NormalTok{),}
                \DecValTok{5000} \SpecialCharTok{/}\NormalTok{ ((}\DecValTok{25}\SpecialCharTok{*}\DecValTok{60} \SpecialCharTok{+} \FloatTok{51.97}\NormalTok{) }\SpecialCharTok{{-}}\NormalTok{ (}\DecValTok{13}\SpecialCharTok{*}\DecValTok{60} \SpecialCharTok{+} \FloatTok{25.00}\NormalTok{)),}
                \DecValTok{5000} \SpecialCharTok{/}\NormalTok{ ((}\DecValTok{41}\SpecialCharTok{*}\DecValTok{60} \SpecialCharTok{+} \FloatTok{57.64}\NormalTok{) }\SpecialCharTok{{-}}\NormalTok{ (}\DecValTok{25}\SpecialCharTok{*}\DecValTok{60} \SpecialCharTok{+} \FloatTok{51.97}\NormalTok{)))}

\NormalTok{profile    }\OtherTok{\textless{}{-}} \FunctionTok{c}\NormalTok{(}\AttributeTok{up =} \FloatTok{0.4}\NormalTok{, }\AttributeTok{flat =} \FloatTok{0.3}\NormalTok{, }\AttributeTok{down =} \FloatTok{0.3}\NormalTok{)}
\NormalTok{speed\_mult }\OtherTok{\textless{}{-}} \FunctionTok{c}\NormalTok{(}\AttributeTok{up =} \FloatTok{0.7}\NormalTok{, }\AttributeTok{flat =} \FloatTok{1.0}\NormalTok{, }\AttributeTok{down =} \FloatTok{1.3}\NormalTok{)}

\NormalTok{A\_beard }\OtherTok{\textless{}{-}} \FunctionTok{runif}\NormalTok{(N, }\FloatTok{0.005}\NormalTok{, }\FloatTok{0.015}\NormalTok{)}
\NormalTok{Cd\_body }\OtherTok{\textless{}{-}} \FunctionTok{runif}\NormalTok{(N, }\FloatTok{0.7}\NormalTok{, }\FloatTok{0.9}\NormalTok{)}
\NormalTok{P\_total }\OtherTok{\textless{}{-}} \FunctionTok{rnorm}\NormalTok{(N, }\DecValTok{400}\NormalTok{, }\DecValTok{30}\NormalTok{)}
\NormalTok{alpha   }\OtherTok{\textless{}{-}} \FunctionTok{runif}\NormalTok{(N, }\FloatTok{0.25}\NormalTok{, }\FloatTok{0.40}\NormalTok{)}

\NormalTok{v3\_mean }\OtherTok{\textless{}{-}} \FunctionTok{numeric}\NormalTok{(N)}
\ControlFlowTok{for}\NormalTok{ (k }\ControlFlowTok{in} \DecValTok{1}\SpecialCharTok{:}\DecValTok{3}\NormalTok{) \{}
\NormalTok{  v\_obs }\OtherTok{\textless{}{-}}\NormalTok{ v\_segments[k]}
\NormalTok{  scale }\OtherTok{\textless{}{-}}\NormalTok{ v\_obs }\SpecialCharTok{/} \FunctionTok{sum}\NormalTok{(profile }\SpecialCharTok{*}\NormalTok{ speed\_mult)}
  \ControlFlowTok{for}\NormalTok{ (seg }\ControlFlowTok{in} \FunctionTok{names}\NormalTok{(profile)) \{}
\NormalTok{    v }\OtherTok{\textless{}{-}} \FunctionTok{rnorm}\NormalTok{(N, speed\_mult[seg] }\SpecialCharTok{*}\NormalTok{ scale, }\FloatTok{0.2}\NormalTok{);  v[v }\SpecialCharTok{\textless{}} \DecValTok{1}\NormalTok{] }\OtherTok{\textless{}{-}} \DecValTok{1}
\NormalTok{    v3\_mean }\OtherTok{\textless{}{-}}\NormalTok{ v3\_mean }\SpecialCharTok{+}\NormalTok{ profile[seg] }\SpecialCharTok{*}\NormalTok{ v}\SpecialCharTok{\^{}}\DecValTok{3}
\NormalTok{  \}}
\NormalTok{\}}
\NormalTok{v3\_mean }\OtherTok{\textless{}{-}}\NormalTok{ v3\_mean }\SpecialCharTok{/} \DecValTok{3}

\NormalTok{P\_drag\_body  }\OtherTok{\textless{}{-}} \FloatTok{0.5} \SpecialCharTok{*}\NormalTok{ rho }\SpecialCharTok{*}\NormalTok{ Cd\_body }\SpecialCharTok{*}\NormalTok{ A\_body }\SpecialCharTok{*}\NormalTok{ v3\_mean}
\NormalTok{P\_drag\_beard }\OtherTok{\textless{}{-}}\NormalTok{ P\_drag\_body }\SpecialCharTok{*}\NormalTok{ (A\_beard }\SpecialCharTok{/}\NormalTok{ A\_body) }\SpecialCharTok{*}\NormalTok{ alpha}\SpecialCharTok{\^{}}\DecValTok{3}
\NormalTok{delta\_time   }\OtherTok{\textless{}{-}}\NormalTok{ T\_total }\SpecialCharTok{*}\NormalTok{ (P\_drag\_beard }\SpecialCharTok{/}\NormalTok{ P\_total)}

\NormalTok{p\_active        }\OtherTok{\textless{}{-}} \FunctionTok{runif}\NormalTok{(N, }\FloatTok{0.05}\NormalTok{, }\FloatTok{0.20}\NormalTok{)   }\CommentTok{\# fraction of race in beard{-}exposed posture}
\NormalTok{delta\_time\_real }\OtherTok{\textless{}{-}}\NormalTok{ delta\_time }\SpecialCharTok{*}\NormalTok{ p\_active}
\end{Highlighting}
\end{Shaded}

\begin{center}\rule{0.5\linewidth}{0.5pt}\end{center}

\subsection{The result}\label{the-result}

\pandocbounded{\includegraphics[keepaspectratio]{mieto-wassberg-useR2026_files/figure-pdf/posterior-plot-1.png}}

\textbf{0.01 s is well inside the bulk} of the distribution.

\begin{center}\rule{0.5\linewidth}{0.5pt}\end{center}

\subsection{What the plot says --- and doesn't
say}\label{what-the-plot-says-and-doesnt-say}

\textbf{What it says}

A margin of 0.01 s is \textbf{not anomalous} under physically plausible
drag assumptions.

The folk theory is not absurd. It is \emph{consistent} with the physics.

\textbf{What it does NOT say}

The beard \emph{caused} Mieto to lose.

We have no counterfactual race without the beard. We have a
\emph{distribution} of plausible drag-time costs.

The right Bayesian sentence: \emph{``0.01 s is consistent with this
model.''} The wrong one: \emph{``the beard cost him gold.''}

\begin{center}\rule{0.5\linewidth}{0.5pt}\end{center}

\subsection{The general lesson}\label{the-general-lesson}

This is a \textbf{template} for a very common task in applied data
science:

\begin{itemize}
\tightlist
\item
  You have a small effect.
\item
  Someone tells you it is ``negligible.''
\item
  You have no clean experiment available.
\item
  You have \textbf{physics, plausible ranges, and R}.
\end{itemize}

Forward simulation lets you ask:

\begin{quote}
\emph{Is the observed effect within the distribution of outcomes that
the system can plausibly produce?}
\end{quote}

You don't need a posterior over parameters. You need a posterior over
\textbf{consequences}.

\begin{center}\rule{0.5\linewidth}{0.5pt}\end{center}

\subsection{Where this approach earns its
keep}\label{where-this-approach-earns-its-keep}

\begin{longtable}[]{@{}
  >{\raggedright\arraybackslash}p{(\linewidth - 4\tabcolsep) * \real{0.1356}}
  >{\raggedright\arraybackslash}p{(\linewidth - 4\tabcolsep) * \real{0.3559}}
  >{\raggedright\arraybackslash}p{(\linewidth - 4\tabcolsep) * \real{0.5085}}@{}}
\toprule\noalign{}
\begin{minipage}[b]{\linewidth}\raggedright
Domain
\end{minipage} & \begin{minipage}[b]{\linewidth}\raggedright
``Negligible'' effect
\end{minipage} & \begin{minipage}[b]{\linewidth}\raggedright
Forward simulation question
\end{minipage} \\
\midrule\noalign{}
\endhead
\bottomrule\noalign{}
\endlastfoot
Finance & A 2 bp execution slippage & What return distribution would
erase it? \\
Compliance & A rule threshold tweak & How many flagged cases shift,
plausibly? \\
Manufacturing & A 0.1 °C drift & Yield consequence under realistic
load? \\
Healthcare & A 1 \% adherence change & Population-level outcome
distribution? \\
Sports & 0.01 s & Drag, wax, draft, terrain --- which matters? \\
\end{longtable}

Every row is the \textbf{same R script} with different priors.

\begin{center}\rule{0.5\linewidth}{0.5pt}\end{center}

\subsection{Why R, specifically}\label{why-r-specifically}

\textbf{The model fits on one screen.}

\begin{itemize}
\tightlist
\item
  \texttt{runif}, \texttt{rnorm}, a \texttt{for} loop, \texttt{ggplot2}.
\item
  No probabilistic programming framework needed for the forward pass.
\item
  Reviewers can read it in 30 seconds.
\end{itemize}

\textbf{The plot is the argument.}

\begin{itemize}
\tightlist
\item
  \texttt{geom\_density} + \texttt{geom\_vline} = the entire claim.
\item
  Reproducible from one \texttt{.qmd} file.
\item
  Priors visible, code visible, conclusion visible.
\end{itemize}

\textbf{Transparency is the deliverable.} The number is secondary.

\begin{center}\rule{0.5\linewidth}{0.5pt}\end{center}

\subsection{Honest limitations}\label{honest-limitations}

\begin{itemize}
\tightlist
\item
  \textbf{Course profile is a proxy.} The 1980 Mt. Van Hoevenberg trail
  is approximated by today's terrain mix (40/30/30).
\item
  \textbf{Weather is assumed calm.} Wind would dominate any beard
  effect.
\item
  \textbf{The \texttt{alpha} slipstream factor} is a back-of-envelope
  shielding estimate, not a CFD result.
\item
  \textbf{\texttt{p\_active}} --- the fraction of the race in
  beard-exposed posture --- is the weakest prior.
\item
  \textbf{No correlation structure} between the four uncertain
  parameters.
\end{itemize}

Every one of these is a knob. Every knob is a \texttt{runif(...)} away
from being tested.

\begin{center}\rule{0.5\linewidth}{0.5pt}\end{center}

\subsection{Extensions (homework, if you
like)}\label{extensions-homework-if-you-like}

\begin{itemize}
\tightlist
\item
  \textbf{Sensitivity analysis}: which prior moves
  \(P(\Delta t \geq 0.01\,s)\) most?
\item
  \textbf{Replace \texttt{runif} with empirical distributions} for
  \(C_d\) from wind-tunnel literature.
\item
  \textbf{Add wind.} A 1 m/s headwind dwarfs the beard.
\item
  \textbf{Add Wassberg.} Two skiers, two posteriors, one difference
  distribution.
\item
  \textbf{Hierarchical version}: pool across the four cross-country
  events at Lake Placid.
\end{itemize}

All of this is plain R. None of it requires more than 100 lines.

\begin{center}\rule{0.5\linewidth}{0.5pt}\end{center}

\subsection{Takeaways}\label{takeaways}

\begin{enumerate}
\def\labelenumi{\arabic{enumi}.}
\item
  \textbf{Forward simulation is underused.} When data is sparse but
  physics is known, simulate first, fit later.
\item
  \textbf{A prior is a hypothesis.} Wide priors are not weakness ---
  they are an honest test of robustness.
\item
  \textbf{The plot is the argument.} A single density with a single
  dashed line settled a 46-year-old question.
\end{enumerate}

\emph{A negligible margin is only negligible relative to a model.}

\begin{center}\rule{0.5\linewidth}{0.5pt}\end{center}

\subsection{}\label{section-1}

Thank you.

Kristian Vepsäläinen

M.Sc. Mathematics · M.Eng. Cyber Security Data scientist · open-data \&
simulation specialist

🌐 kristianvepsalainen.com ✉️ kristian.vepsalainen@proton.me

Code \& slides: github.com/kristianvepsalainen/useR2026 ~·~ CC BY 4.0




\end{document}
